% Mechanically converted from the audited Markdown manuscript.
\documentclass[11pt]{article}
\usepackage[a4paper,margin=27mm]{geometry}
\usepackage[T1]{fontenc}
\usepackage[utf8]{inputenc}
\usepackage{lmodern}
\usepackage{microtype}
\usepackage{amsmath,amssymb,amsthm,mathtools}
\usepackage{xcolor}
\usepackage[colorlinks=true,linkcolor=blue!55!black,citecolor=blue!55!black,urlcolor=blue!65!black]{hyperref}
\usepackage{xurl}
\usepackage{fancyhdr}

\newtheorem{theorem}{Theorem}
\numberwithin{equation}{section}
\pagestyle{fancy}
\fancyhf{}
\fancyhead[L]{Carptopus}
\fancyhead[R]{Connected-set polynomial counterexamples}
\fancyfoot[C]{\thepage}
\setlength{\headheight}{14pt}
\setlength{\parindent}{0pt}
\setlength{\parskip}{0.42em}
\setlength{\emergencystretch}{4em}

\title{Counterexamples to unimodality\\of the connected set polynomial}
\author{Carptopus\\\href{mailto:carptopus@163.com}{carptopus@163.com}}
\date{Manuscript, 2 September 2026}

\hypersetup{
  pdftitle={Counterexamples to unimodality of the connected set polynomial},
  pdfauthor={Carptopus},
  pdfsubject={Infinite counterexample families and the complete minimum-order classification},
  pdfkeywords={connected set polynomial, unimodality, graph polynomial, connected induced subgraph, counterexample, graph enumeration}
}

\begin{document}
\maketitle

\begin{abstract}
For a finite graph $G$, let $C(G;x)=\sum_{k\geq1}c_k(G)x^k$, where $c_k(G)$ counts the
$k$-vertex subsets inducing a connected subgraph. A question of Mol asks whether $C(G;x)$ is
unimodal whenever $G$ is connected and $|E(G)|\geq|V(G)|$. For every order $n\geq8$, we give
four explicit pairwise nonisomorphic counterexamples, with edge counts
$n+1,n+2,n+2,n+3$. Their polynomials follow from a common rooted-block decomposition and exhibit
the fixed local pattern $10,11,10$. We also enumerate all connected graphs of order eight
satisfying the edge condition and prove that exactly four isomorphism classes are counterexamples;
they are precisely the first members of our four families. No eligible graph of smaller order is a
counterexample. Thus the minimum counterexample order is eight.
\end{abstract}

\section{Introduction}

The connected set polynomial of a finite simple graph $G$ is
\begin{equation}
C(G;x)=\sum_{k=1}^{|V(G)|}c_k(G)x^k,
\end{equation}
where $c_k(G)$ is the number of $k$-element sets whose induced subgraph is connected. This
polynomial is the one-component part of the induced-subgraph component polynomial introduced by
Tittmann, Averbouch and Makowsky~\cite{tittmann}, and its coefficients are closely related to node
reliability.

Mol asked whether the coefficient sequence is always unimodal for a connected $n$-vertex graph
with at least $n$ edges; the same paragraph reports a computational check through order eight.
The statement appears in the conclusion chapter, under ``Connected Set Polynomial,'' of the 2016
thesis cited below~\cite{mol}. The theorem below answers the question negatively and determines the
complete minimum-order boundary.

\begin{theorem}\label{thm:main}
For every $n\geq8$, there are four pairwise nonisomorphic connected $n$-vertex graphs with at
least $n$ edges whose connected set polynomials are not unimodal. Their edge counts are
\[
n+1,\qquad n+2,\qquad n+2,\qquad n+3.
\]
Among all connected eight-vertex graphs with at least eight edges, exactly four isomorphism classes
have a nonunimodal connected set polynomial, and they are the four order-eight graphs in these
families. No eligible graph of order at most seven is a counterexample.
\end{theorem}

\section{Four rooted-block families}

Fix $d\geq2$. Join rooted graphs $L$ and $R$ by a path of length $d$, identifying its endpoints
with the roots $u\in L$ and $v\in R$.

For the left block choose either $D=K_4-e$, rooted at a vertex of degree three, or $K=K_4$. For
the right block choose either $S=K_{1,2}$, rooted at its centre, or $T=K_3$. Denote the four
resulting graphs by
\[
G_d^{D,S},\qquad G_d^{D,T},\qquad G_d^{K,S},\qquad G_d^{K,T}.
\]
Let $\varepsilon$ be one when the left block is $K$ and zero otherwise, and let $\eta$ be one when
the right block is $T$ and zero otherwise. Every graph has
\[
n=d+6,\qquad m=n+1+\varepsilon+\eta.
\]
The four graphs are pairwise nonisomorphic. Their block decompositions recover whether the left
nontrivial block is $K_4-e$ or $K_4$, and whether the right endpoint lies in a triangle.

\section{Rooted decomposition}

For a rooted graph $(H,r)$, let $A_H(x)$ count the nonempty connected induced vertex sets
containing $r$. In both left blocks the root is adjacent to all three other vertices, so
\[
A_L(x)=x(1+x)^3.
\]
Similarly, both right blocks satisfy
\[
A_R(x)=x(1+x)^2.
\]
The sum of the two ordinary connected set polynomials is
\[
C(L;x)+C(R;x)=7x+(7+\varepsilon+\eta)x^2+5x^3+x^4.
\]

We partition a connected induced set according to whether it lies in an endpoint block, lies wholly
inside the $d-1$ internal path vertices, meets one endpoint block and an initial path segment, or
meets both endpoint blocks. These classes are disjoint and exhaustive. Consequently,
\begin{align}
C(G_d;x)={}&7x+(7+\varepsilon+\eta)x^2+5x^3+x^4
 +\sum_{k=1}^{d-1}(d-k)x^k \notag\\
&+(A_L(x)+A_R(x))(x+x^2+\cdots+x^{d-1})
 +A_L(x)A_R(x)x^{d-1}.
\tag{3.1}
\end{align}
Expanding (3.1) yields
\begin{align}
C(G_d;x)={}&nx+(n+1+\varepsilon+\eta)x^2+(n+3)x^3 \notag\\
&+\sum_{k=4}^{d+2}(n+6-k)x^k
 +11x^{d+3}+10x^{d+4}+5x^{d+5}+x^{d+6}.
\tag{3.2}
\end{align}
The sequence decreases after $c_3=n+3$, reaches $c_{d+2}=10$, and then rises to
$c_{d+3}=11$. Hence all four polynomials are nonunimodal. Since $d=n-6$, this proves the
infinite-family part of Theorem~\ref{thm:main}.

\section{The minimum-order classification}

At $d=2$, equation (3.2) becomes
\[
(8,\ 9+\varepsilon+\eta,\ 11,\ 10,\ 11,\ 10,\ 5,\ 1).
\]
We checked Brendan McKay's complete \texttt{graph8c.g6} catalogue of 11,117 nonisomorphic connected
graphs of order eight. Exactly 11,094 have at least eight edges, and exactly four of these have a
nonunimodal connected set polynomial. This is a complete catalogue scan, not a generated sample.
Their graph6 encodings are
\begin{flushleft}
\ttfamily
G?B@vG\\
G?\textasciigrave{}crg\\
G?aJeW\\
GCQRTg
\end{flushleft}
Isomorphism tests identify them with the four $d=2$ constructions above. An independent scan of
the NetworkX graph atlas finds no eligible counterexample of order at most seven. The connected-set
coefficients were computed by two independent routines: direct subset connectivity and boundary
expansion. They agree on every tested graph. This proves the finite classification in
Theorem~\ref{thm:main}.

The catalogue is available from
\url{https://users.cecs.anu.edu.au/~bdm/data/graph8c.g6}. Its expected SHA-256 digest is
\begin{center}
\small\ttfamily
0002354F1AB3344A2706626A037AD15367BF23A2163AA68F552C3A169CA9A036.
\end{center}
The verification used Python 3 and NetworkX 3.6.1. The supplied verification directory records the
complete command and the logical boundary between the formula checks, the atlas scan and the full
order-eight catalogue scan.

\section{Scope}

The theorem does not classify all graphs with nonunimodal connected set polynomials. In particular,
we do not claim that four is the number of counterexamples in orders greater than eight, or that no
counterexample with exactly $n$ edges exists. The result identifies four explicit infinite
families and completely settles only the minimum-order boundary.

{\small
\begin{thebibliography}{9}

\bibitem{dossou}
A.~Dossou-Olory,
\emph{Graphs and unicyclic graphs with extremal number of connected induced subgraphs},
arXiv:1812.02422.

\bibitem{mol}
L.~Mol,
\emph{On Connectedness and Graph Polynomials},
PhD thesis, Dalhousie University, 2016, conclusion chapter, ``Connected Set Polynomial'' subsection;
\href{https://central.bac-lac.gc.ca/.item?app=Library&id=TC-NSHD-71408&oclc_number=1033184059&op=pdf}{stable thesis record and PDF}.

\bibitem{tittmann}
P.~Tittmann, I.~Averbouch and J.~A.~Makowsky,
\emph{The Enumeration of Vertex Induced Subgraphs with respect to the Number of Components},
European Journal of Combinatorics 32 (2011), arXiv:0812.4147.

\end{thebibliography}
}

\paragraph{AI-assisted research and writing disclosure.}
Carptopus is the author responsible for the mathematical claims and the public version of this
manuscript. OpenAI Codex was used as an AI-assisted tool for literature organization, proof
development and stress testing, exact verification, adversarial auditing, and manuscript
preparation.

\end{document}

